\documentclass[12pt]{article}
\usepackage[utf8]{inputenc}
\usepackage[T1]{fontenc}
\usepackage{hyperref}
\title{Tauade Language Resources for K-12 Curriculum Development}
\author{Compiled by Claude Sonnet 5 (claude-sonnet-5) for Andr\'as Kornai}
\date{}
\begin{document}
\maketitle

\section{Introduction}
Tauade (ISO 639-3 ttd) is a Trans-New Guinea language of Papua New Guinea. Estimates run from about 7,000 (2000 census, per Wikipedia) to 12,000--13,000 speakers (other sources). Little material was found online.

\section{Primary Educational Resources}
None found.

\section{Linguistic and Reference Materials}
\begin{itemize}
\item Tauade--English dictionary by Robert and Verna Stutzman, published on Webonary (SIL International, 2016). \url{https://www.webonary.work/tauade/?p=146817}
\item Wikipedia: \url{https://en.wikipedia.org/wiki/Tauade_language}
\end{itemize}

\section{Religious and Community Resources}
Joshua Project lists only Bible portions (dated 1991) and no New Testament. It also reports that a translation team worked at Sopu from 1987 to 1991 and that the project ended after an attack on team members. \url{https://joshuaproject.net/languages/ttd}

\section{Digital and Audio Resources}
No audio or video resources were found. The Webonary dictionary is the only digital resource found.

\section{Curriculum Development Context}
No school use was found, which matches the reference description. The Webonary dictionary is the one usable base for vocabulary work. Any further material would need to come from fieldwork or from the dictionary compilers' wider documentation.
\end{document}